% !TeX root = ../papers/nc-p3.tex
\section{七余因子の中間退化の全特殊値と余次元五の排除}
\label{sec:trisecant-middle-degeneracy}
\begin{theorem}\label{thm:trisecant-middle-degeneracy}
四辺有限窓の整な七余因子の場合、標準中間座標の二十点上で六行四列の評価行列は必ず階数3である。
従って非被約の場合も二つの三次因子と標準中間座標の三スキームは同型で、
両隣の標準辺を復元できる。相異なる四次曲線が交わらないパラメータ開集合内で、
中間に階数2点を持つパラメータは余次元5の既約閉集合をなし、その全体が排除される。
\end{theorem}
\begin{proof}[次数三の相対Koszul余核]
中間点を $c=[l:m]$ とし $M_c(x)=l T(sx)+m T(tx)$ を $S_3\to S_5$ の写像とする。
その最大行列式スキーム $Z$ は既出のアフィン性で有限長20であり、全台点で $l,m$ は独立である。
$X$ の点環全体で中間座標が復元される。実際 $T(sx),T(tx)$ は全ての台点で独立で、
階数1のBezout行列は共通三次因子を与える。この相対表示は
$\bigwedge^2 S_4\simeq\operatorname{Sym}^2 S_3$ で $y^2$ の逆が $sy\wedge ty$ となる恒等式からも分かる。
単元小行列式で三次因子を持つ二基底の変更を解けば $[l:m]$ はスキーム上で一意である。
従って $X$ は核直線を印した入射に入力射影を逆として閉埋込みされる。

$Z$ の周囲で $M:\mathcal O(-1)^4\to\mathcal O^6$ の転置による線型方程式の
内部次数3のKoszul複体を取る。その自由項は右から
\[
              \mathcal O(3)^{20},\quad \mathcal O(2)^{60},\quad
              \mathcal O(1)^{60},\quad \mathcal O^{20}.
\]
$Z$ の外では転置が全射で線型関係は二次元、そのKoszulホモロジーは内部次数2までなので
この三次部分は完全である。三次元正則な底で長さ3、ホモロジーが有限台なので終端余核を除いて完全となる。
余核 $U$ は線型座標加群の対称三乗であり、Euler数から長さ20である。
固定 $c$ のファイバーで入力は核の射影線型空間に入り、入力三次式は $X$ 上の全切断に全射なので
各ファイバーにも全射である。中山の補題から $U$ は $X$ 上の相対三次切断の押出しへ全射する。
両長さが20なので同型である。階数2の $c$ なら $U$ の閉繊維は二変数三次の次元4であるが、
入力直線切断の長さは3以下である。階数1でも同じく平面三次の10に対し平面切断長は7以下である。
階数0なら固定 $l,m$ により全 $x$ で $T(sx)$ は同じ一次因子を持ち、二つの四次曲線に共通点が生じる。
よって全階数は3で、核入射は $Z$ 自体、部分スキーム $X$ と等長なので一致する。
\end{proof}

\begin{proof}[排除されたパラメータの実在と次元]
中間空間 $\mathrm{PGL}_2$ の一点を固定し、二項変数の変更で
$M=\partial T:t T(sx)-s T(tx)$ と書く。線型写像
\[
        \partial:\operatorname{End}S_4\to\operatorname{Hom}(S_3,S_5)
\]
の核はスカラーだけである。零なら任意の一次式 $v$ の空間 $vS_3$ を $T$ が保つから、
双対では有理正規四次曲線の全評価直線を保つ。これらは一つの射影枠を含むのでスカラーとなる
（非可逆でも、既知のスカラーを加えて可逆にしてよい）。次元差1なので全射である。
固定中間で階数2以下となる条件は従って既約な余次元 $(6-2)(4-2)=8$ の決定条件である。
曲線が交わらなければ階数1以下はどの中間でも不可である。実際像の五次式が共通根を持つと
その根での $T^*$ の評価は所定の四次評価に比例し、共通点を生じる。

この既約軌跡に二曲線が交わらず、行列式八次式も整な点が存在する。
$M(x)=F(z)\alpha(x)+G(z)\beta(x)$ とし、二つの五次式は基点なし、
$\alpha,\beta$ の直線は双対三次有理正規曲線を避けるように取る。
$\partial T=M$ の解にスカラーを加えて可逆にする。もし $T^*\operatorname{ev}_z$ が一点評価なら、
$M(-)(z)$ は零または三次の一点評価である。しかしこの汎関数は非零で指定直線に入り矛盾する。
八次式の母線 $\ell_z$ は平面 $F(z)\alpha+G(z)\beta=0$ に入る。
$z\mapsto\ell_z$ のPlücker次数は8である。二つの四次評価行は全点で独立なのでその外積に基点がない。
一般の母線は軸 $\alpha=\beta=0$ そのものではなく、その平面を一意に定める。
平面の射の次数は5なので、母線写像の像への次数は $\gcd(5,8)=1$ を割る。
一般の平面内の五本の母線は従って相異なり、軸外の一般入力は一意な母線に入る。
母線族は結果式の零点全てを覆い、そのパラメータ線束のPlücker次数8と一意性から像は次数8の整曲面である。

全中間で印を動かした入射は $\mathrm{PGL}_2$ の推移作用により既約である。
中間の退化は射影的で常に有限だから、その開パラメータへの像は閉で、有限射の次元で余次元は $8-3=5$ となる。
なお階数3の印 $c$ では、$T(sx)=s y,T(tx)=t y$ と $c$ を規格化すれば、中間点の接空間次元は
$\dim(y S_2\cap\operatorname{im}\partial T)$ である。行列式の接写像は $\ker M_c$ から余核への微分であり、
三つの中間接方向による $M'_c(x)$ の包がちょうど $y S_2$ だからである。
\end{proof}

\path{computations/verify_trisecant_middle_degeneracy.py} は $\partial$ の階数24、
基点なし両端五次単項式と中間二係数による階数2の可逆有理標本、三次余核の長さ20を検算する。
